% Generated draft: requires mathematical and source review.
\documentclass[11pt]{article}
\usepackage[margin=25mm]{geometry}
\usepackage{fontspec}
\setmainfont{Latin Modern Roman}
\usepackage{amsmath,amssymb,mathtools,unicode-math}
\setmathfont{Latin Modern Math}
\usepackage{enumitem,needspace,xurl,hyperref}
\setlength{\parindent}{0pt}
\setlength{\parskip}{5pt}
\setlength{\emergencystretch}{3em}
\newcommand{\N}{\mathbb{N}}
\newcommand{\Z}{\mathbb{Z}}
\newcommand{\Q}{\mathbb{Q}}
\newcommand{\R}{\mathbb{R}}
\newcommand{\C}{\mathbb{C}}
\newcommand{\F}{\mathbb{F}}
\newcommand{\A}{\mathbb{A}}
\newcommand{\PP}{\mathbb{P}}
\newcommand{\E}{\mathbb{E}}
\newcommand{\eps}{\varepsilon}
\newcommand{\dd}{\,\mathrm{d}}
\begin{document}
\section*{UnsolvedMath ID: 11000134}
\textbf{Source problem code:} catalogue entry
\subsection{Definitions and mathematical statement}
Fix the surface $F$ and let $\mathcal P_4(F)$ be the mapping classes produced by the Dehn-twist recipe in the source's Theorem 4. That theorem and its recipe are not supplied. A pseudo-Anosov (pA) class has stretch factor $\lambda(f)>1$ and invariant transverse measured foliations scaled by reciprocal factors.

Problem 4 asks whether
\[
 \forall f\in\operatorname{Mod}(F)\text{ pseudo-Anosov},\quad
 \exists n\in\Z_{\geq1}:\quad f^n\in\mathcal P_4(F).
\]
Thus ``virtually all'' means that every pA class has an iterate produced by this particular recipe. It does not mean a density-one assertion, passage to a finite cover of the surface, or merely using some other construction. The registry relates the question to Penner's construction, but this label alone does not supply all the omitted hypotheses and allowed twist words.

The source traces the question to [23,25] and its relationship to the Ehrenpreis Conjecture, identified there as Problem 14, to [10]. It also cites Mosher [42] and joint work of the chapter author with Papadopoulos [45] for other descriptions of pA maps up to iteration; these are described as combinatorial rather than Dehn-twist descriptions and therefore do not automatically answer this question.

The following extracted spectrum discussion defines
\[
 \Sigma(F)=\{\log\lambda(f):f\text{ is pseudo-Anosov on }F\}\subset\R_{>0}.
\]
It identifies this with the Teichm\"uller geodesic length spectrum of $\mathcal M(F)$ and reports discreteness, explaining that dilatations occur as spectral radii of integral-linear Perron--Frobenius symplectomorphisms in dimensions bounded in terms of the topology of $F$. These are supplied background assertions, without the dimensional bounds or surface conventions needed to restate their exact theorems.
\subsection{Short English statement}
Does every pseudo-Anosov mapping class have a positive power arising from the particular Dehn-twist recipe of the omitted Theorem 4? Other combinatorial constructions up to iteration do not answer that exact recipe question.
\subsection{Sources}
\begin{itemize}
\item Local registry: \url{data/unsolvedmath/problems.json}, permanent UnsolvedMath ID $11000134$, exact source code \texttt{catalogue entry}.
\item Farb, \emph{Problems on Mapping Class Groups and Related Topics}, Problem 4, PDF page 117: \url{https://www.math.uchicago.edu/~farb/papers/mcgbook.pdf}. Supplied access date: 2026-07-29; extraction: pdf-font-structure; repeated local numbering disambiguated by page among 329 headings.
\item Supplied citations: [23,25] for the recipe question, [10] for the relation to Ehrenpreis/Problem 14, Mosher [42], and the chapter author with Papadopoulos [45]. Theorem 4, Problem 14 and the bibliography entries are absent. The registry calls this Penner-related without a precise modern reference.
\end{itemize}
\end{document}
